% ============================================================================
% FIRMWARE CHAINING, VALUE EMISSION, AND PROGRAM COMPOSITION
% ============================================================================
\subsection{Firmware Chaining and Value Emission}
\label{sec:firmware}

The program region of a single Value holds 98 instruction slots.  This
is not, however, the limit of a program's length.  Because the firmware
index register (\texttt{fw}) and the program counter (\texttt{pc}) are
both addressable by the same boolean operations used for all other
computation, a program can \emph{chain} to the next firmware by writing
a new firmware index into \texttt{fw} and resetting \texttt{pc} to zero
as its final two instructions.  On the next execution cycle, the
bootloader loads the indicated firmware and execution continues from the
beginning of the new program.  An arbitrarily long computation can
therefore be expressed as a sequence of firmware segments, each fitting
within 98 slots, linked end-to-end through the \texttt{fw}/\texttt{pc}
handoff.

\subsubsection{Firmware Compilation}

Each firmware segment is specified as a sequence of textual
instructions in the format \texttt{source~destination~operation}.
Operands may name a register (\texttt{r0}--\texttt{r5}, \texttt{pc},
\texttt{fw}), a pointer dereference, or a fourteen-bit immediate
value.  The operation field is a four-bit binary literal selecting one
of the sixteen truth tables.  A compiler translates each line into a
packed thirty-two-bit instruction word.  At runtime, any compiled
segment can be loaded into a Value's program region by setting
\texttt{fw} to the corresponding index.

\subsubsection{The Firmware Library}

Six firmware segments are defined.  Each performs a focused structural
operation; together, they compose into longer pipelines through
chaining.

The \textbf{bootloader} initializes a frame's registers to known
values and prepares it for subsequent execution.  Every firmware chain
begins here: a chain is expressed as bootloader~$\to$~segment$_1$~$\to
$~segment$_2$~$\to \cdots$, with each segment ending by writing the
next index into \texttt{fw} and resetting \texttt{pc}.

The \textbf{affinity} program performs two operations.  First, it
accumulates a structural summary of the executing frame's data region
through a running boolean combination across the token words.  Second,
it adopts linking metadata from a partner frame, establishing an edge
in the implicit graph that connects related Values.  Crucially, the
result of affinity is not merely a modification of the executing frame:
it is intended to \emph{emit} a new Value whose data region encodes
the folded structure---the residual after shared components between
the two participating frames have been identified and factored out.
This emission deposits the new Value back into the Region, where it
becomes available for further interaction.

The \textbf{viral} program copies the executing frame's program region
into a partner frame's program region.  Its final instructions then
write a specified firmware index into the partner's \texttt{fw}
register and reset the partner's \texttt{pc}, so that the partner
wakes up running the next segment in the chain.  A typical deployment
is viral~$\to$~affinity: the viral step propagates the query logic to
a neighboring frame, and the affinity step evaluates structural overlap
between the query and the neighbor's data.

The \textbf{shatter} program applies a three-pass AND operation that
decomposes a frame's data content, progressively isolating subsets of
active bits into distinct register states.  The \textbf{tombstone}
program zeroes structural metadata---linking pointers, state fields,
and the affinity bitmask---to mark a frame as inactive.  The
\textbf{wipe} program zeroes the data region while preserving
structural metadata.

\subsubsection{Graph Construction Through Folding}

When the affinity firmware processes a set of Values that share a
common component---for example, the predicate ``is in the'' appearing
in three distinct propositions---the shared component factors out and
becomes a structural node in the resulting graph.  The mechanism is
boolean cancellation: applying XNOR across two data regions yields all
ones where the regions match, and the residual (obtained via AND-NOT
or XOR) isolates the distinguishing content.

Consider three propositions: ``Sandra is in the Garden,'' ``Roy is in
the Kitchen,'' ``Harold is in the Kitchen.''  After affinity processing,
the shared predicate collapses into a single label node that links to
the set of subjects $\{$Sandra, Roy, Harold$\}$, while each subject
links to its associated location.  The link pointers (word~65) and
affinity bitmasks (word~64) encode these edges.  The emitted Values
representing the folded structure are deposited back into the Region.

A query such as ``Where is Roy?'' then traverses this graph: the
predicate component cancels against the label node via XNOR, the
subject component cancels against the Roy node, and the remaining edge
resolves to ``Kitchen.''  The traversal cost is proportional to the
graph depth, not the total number of stored Values.

\subsubsection{Firmware Chains as Pipelines}

The chaining mechanism allows firmware segments to compose into
multi-stage pipelines.  A retrieval pipeline, for instance, might
proceed as:

\begin{enumerate}[nosep]
  \item \textbf{Viral:} propagate the query program to a neighbor.
  \item \textbf{Affinity:} evaluate structural overlap; emit a new
        Value encoding the residual.
  \item \textbf{Shatter:} decompose the residual to isolate individual
        features.
  \item \textbf{Affinity} (second pass): match decomposed features
        against the neighbor's link targets.
\end{enumerate}

Each transition is a two-instruction \texttt{fw}/\texttt{pc} handoff.
The pipeline length is limited only by the number of defined firmware
segments, not by the 98-slot capacity of a single program region.

\subsubsection{Self-Modification and Operator Discovery}

Because the program region is part of the frame's addressable memory,
boolean operations applied during Region mixing can produce instruction
sequences that differ from any of the six firmware templates.  If such
a derived program produces a useful transformation, the viral firmware
can propagate it to neighboring frames.

The selection signal for ``useful'' need not be external.  A frame's
own state region tracks execution outcomes: the TTL counter (word~70)
decrements with each hop, and the state accumulator (word~63) records
whether interactions produced non-trivial residuals.  A program that
leads to successful structural matches---non-zero XNOR overlap with
query targets---will tend to leave its host frame in a state that
other frames recognize as high-affinity, increasing the probability
that the viral mechanism copies it.

This creates a feedback loop: firmware seeds initial behavior; Region
mixing produces novel instruction sequences; the viral mechanism
propagates sequences from frames with favorable state; and the
cycle repeats.  No gradient computation is involved.  The process is
combinatorial, driven by the stochastic mixing order of the Region
and the deterministic execution of boolean instructions.
